\chapter{Программная реализация высокопроизводительного ядра навигации}
\label{chap:implementation}

В данной главе детально описывается архитектура программного ядра маршрутизации, разработанного на языке программирования C++. Рассматриваются процессы трансформации сырых картографических данных в бинарные структуры компактного формата, реализация алгоритма динамического поиска кратчайших путей с применением векторных инструкций, а также интеграция модулей телеметрии и принятия решений.

\section{Архитектура программного комплекса \texttt{traffic\_core}}

Разработанный программный комплекс представляет собой высокопроизводительный Linux-демон, написанный на стандарте C++26. В отличие от классических монолитных транспортных симуляторов (например, SUMO или MATSim), архитектура \texttt{traffic\_core} базируется на принципе строгой изоляции вычислительных контуров и отказе от сильной связности (Tight Coupling).

\subsection{Функциональная декомпозиция модулей}

Для предотвращения деградации кодовой базы и появления монолитных классов («God Objects»), система физически разделена на изолированные статические библиотеки (модули). Каждый модуль имеет собственный \texttt{CMakeLists.txt} и инкапсулирует свою зону ответственности:

\begin{itemize}
    \item \textbf{\texttt{router}} (Навигационное ядро) — чистый \textit{stateless}-модуль In-Memory вычислений. Отвечает исключительно за алгоритмический обход графа (TD-ALT), поддержку 4-арных куч (Heap) и работу с плоскими массивами формата CSR. Не содержит логики симуляции.
    \item \textbf{\texttt{data\_provider}} (Среда и Симулятор) — модуль, отвечающий за 1D-кинематику агентов (модель жидкости — Fluid Model), продвижение физического времени и трансляцию сырых метрик загруженности дорог.
    \item \textbf{\texttt{decision\_engine}} (Модуль принятия решений, МПР) — интеллектуальный «мозг» системы. Анализирует потоки данных от \texttt{data\_provider}, выявляет зарождающиеся заторы с помощью пространственных индексов и отправляет асинхронные запросы (триггеры) в \texttt{router} для локального или глобального перестроения маршрутов.
    \item \textbf{\texttt{core}} — базовый слой инфраструктуры, предоставляющий пулы потоков (Thread Pool), Lock-Free очереди (на базе \texttt{moodycamel::ConcurrentQueue}) и диспетчеризацию задач.
    \item \textbf{\texttt{daemon}} — точка входа (\texttt{main.cpp}), реализующая паттерн \textit{Composition Root}.
\end{itemize}

\subsection{Inversion of Control и паттерн Composition Root}

Связывание независимых подсистем осуществляется через паттерн внедрения зависимостей (Dependency Injection). Модули общаются исключительно через абстрактные контракты (например, интерфейсы \texttt{IAgentDataProvider} или \texttt{ITelemetryPublisher}), ничего не зная о внутренней реализации друг друга. 

Сборка итогового конвейера симуляции происходит в единственной точке приложения — функции \texttt{main()}, выступающей в роли корня композиции (Composition Root). Оркестратор инициализирует объекты в строгой последовательности от ядра к периферии, формируя следующий цикл (Main Loop):
\begin{enumerate}
    \item \textbf{Шаг симуляции:} \texttt{data\_provider} продвигает физику на $\Delta t$ и формирует массивы изменений (дельт).
    \item \textbf{Шаг телеметрии:} Извлеченные легковесные дельты передаются в сетевой слой ZeroMQ для асинхронной отправки на веб-клиенты.
    \item \textbf{Шаг аналитики:} МПР забирает сгенерированные симулятором события-алерты (например, критическое отклонение \texttt{ETA}).
    \item \textbf{Шаг маршрутизации:} МПР запрашивает у пула потоков альтернативные маршруты и возвращает их агентам.
\end{enumerate}

Данный подход превращает архитектуру в легко расширяемый конструктор, позволяющий в будущем (например, при интеграции машинного обучения) прозрачно подменить математический МПР на нейросетевой инференс без переписывания транспортного ядра.

\subsection{Zero-Cost OOP и Data-Oriented Design}

Классический объектно-ориентированный дизайн (с обильным использованием интерфейсов и виртуальных функций) провоцирует промахи кэша инструкций процессора и блокирует агрессивный инлайнинг (inlining) на этапе компиляции. В горячем цикле маршрутизатора это приводит к недопустимому нарушению бюджета времени в 190 наносекунд на одну итерацию обхода узла.

Для сочетания чистой архитектуры с аппаратной скоростью выполнения применена парадигма абстракций с нулевой стоимостью (Zero-Cost OOP):
\begin{enumerate}
    \item \textbf{Запрет на виртуальные вызовы (vtable):} В методах обхода графа и симуляции физики (\texttt{tick}, \texttt{step}, \texttt{calculate\_path}) полностью исключено динамическое связывание.
    \item \textbf{Инкапсуляция плоских данных:} Тяжелые массивы \texttt{row\_ptr} и \texttt{col\_ind} (формат CSR) инкапсулируются внутри классов как приватные поля \texttt{std::vector}. Внешний интерфейс предоставляет безопасные геттеры, помеченные директивами \texttt{inline} и \texttt{[[nodiscard]]}.
    \item \textbf{Compile-Time полиморфизм:} Компилятор LLVM/GCC полностью вырезает вызов таких методов, заменяя их прямым ассемблерным смещением в ОЗУ. Это гарантирует скорость языка ассемблера при сохранении строгой типизации и безопасности памяти современного C++.
\end{enumerate}

Архитектурная схема взаимодействия модулей представлена в Приложении А (рис. \ref{fig:system_arch}).

\section{Офлайн-препроцессинг и сжатие графа}

Для обеспечения микросекундной задержки при маршрутизации вычислительное ядро полностью освобождено от задач парсинга геометрии. Все ресурсоемкие операции вынесены в специализированный офлайн-контур подготовки графа.

Инверсия топологии и реберный граф.
Классический узловой граф (где перекрестки — узлы, дороги — ребра) непригоден для мезоскопической модели, так как при такой архитектуре стоимость маневра (например, запрет поворота) зависит от истории предыдущего шага, что требует внедрения условных переходов в горячем цикле алгоритма. 

В разработанном контуре применена инверсия — направленный реберный граф:
\begin{itemize}
    \item Физические участки дорог (сегменты) становятся узлами нового вычислительного графа.
    \item Маневры на перекрестках (переходы между сегментами) становятся ребрами.
\end{itemize}
Это позволяет жестко внедрить все геометрические штрафы (за торможение, прохождение дуги поворота и разгон) в статический вес связующего маневра. Для графа размерности Московской агломерации генерация связей на перекрестках создает плотную сеть из 1.5 млн возможных маневров.

Хранение в формате компактных разреженных строк и топологическая пересортировка.
Классическое объектное представление графа в куче вызывает массовые промахи кэша процессора из-за случайного доступа к памяти. Для преодоления этого барьера («стены памяти») вся топология сериализуется в формат компактного хранения разреженных строк (CSR). Граф описывается набором плоских массивов (структура массивов):
\begin{enumerate}
    \item массив смещений указателей строк — массив из N+1 элементов, хранящий индексы начала списка смежных маневров для каждого узла.
    \item массив индексов столбцов — массив из E элементов, хранящий идентификаторы целевых сегментов.
    \item массив статических весов — массив из E элементов, хранящий базовое время проезда свободного ребра.
\end{enumerate}

Схематичная структура хранения компактного представления графа в бинарном дампе памяти приведена на рис. \ref{fig:binary_dump}.

\begin{figure}[htbp]
    \centering
    \includesvg[width=0.95\linewidth,height=0.45\textheight,keepaspectratio,inkscapelatex=false]{assets/pdf/binary_dump.drawio.svg}
    \caption{Структура бинарного представления графа компактного формата в памяти}
    \label{fig:binary_dump}
\end{figure}

Для максимизации процента попаданий в кэш процессора первого и второго уровня на этапе сборки дампа применяется алгоритм топологической пересортировки графа. Рекурсивная бисекция графа переназначает идентификаторы узлов таким образом, чтобы топологически связные сегменты получали смежные адреса в памяти. При загрузке узла контроллером памяти (блоками по 64 байта) аппаратно выполняется предвыборка его фактических соседей.

Выравнивание структур данных и предрасчет функций задержки.
В то время как статика графа хранится в формате структуры массивов, динамическое состояние (трафик) меняется непрерывно. Если хранить параметры затора разрозненно, вычисление функции задержки потребует нескольких промахов кэша на один узел. Для решения этой проблемы применена архитектура массива структур (AoS).

Состояние каждого физического сегмента инкапсулируется в структуру состояния динамического узла, которая принудительно выровнена под размер кэш-линии первого уровня директивой выравнивания по границе 64 байт:
\begin{plaincode}{dynamic_node_state}{Структура динамического состояния узла с выравниванием по кэш-линии}
struct alignas(64) DynamicNodeState {
    uint32_t K_magic;                         // Целочисленный множитель (4 байта)
    std::atomic<uint16_t> mpr_penalty;        // Внешний штраф (2 байта)
    uint16_t _reserved;                       // Заполнитель (2 байта)
    std::atomic<uint16_t> volume_buckets[24]; // Корзинки прогнозов (48 байт)
    // Оставшиеся 8 байт (защита от ложного разделения данных)
};
\end{plaincode}

Важнейшим элементом оптимизации является предрассчитанный целочисленный множитель. Для исключения операций деления и возведения в произвольную степень в горячем цикле роутера, базовые константы функции задержки (базовая скорость, емкость сегмента) аналитически сжимаются в единое целое число на этапе сборки графа:
\begin{equation}
    K_{magic} = \left\lfloor \frac{T_{free} \cdot \left( \frac{V_{free}}{5} - 1 \right)}{C_{300}^2} \cdot 2^{20} \right\rfloor
\end{equation}

Умножение на $2^{20}$ позволяет в реальном времени использовать аппаратную арифметику с фиксированной точкой (битовые сдвиги вправо на 20 разрядов) вместо медленных вычислений с плавающей запятой, сокращая время расчета динамического штрафа до 3–4 тактов процессора.

Генерация координационных матриц опорных вершин.
В процессе предварительной обработки также выполняется выборка опорных ориентиров (маяков) для алгоритма динамического поиска по опорным вершинам. Для равномерного покрытия транспортной сети применяется стратегия максимального покрытия в сочетании с алгоритмом граничного минмакса. Офлайн-генератор проводит полные обходы графа алгоритмом Дейкстры, фиксируя 32 наиболее равноудаленных маяка, после чего вычисляет и упаковывает расстояния от каждого узла до всех 32 маяков в бинарный дамп графа. Это формирует базу для вычисления эвристической оценки за константное время на лету.

Конвейер построения вычислительного графа.
Сквозной процесс офлайн-подготовки картографических данных представляет собой строго направленный конвейер (построитель графа), реализующий последовательную цепочку шагов от парсинга сырого картографического файла до выгрузки оптимизированного бинарного дампа графа. 

Основные этапы работы конвейера подготовки графа:
\begin{enumerate}
    \item Импорт геоданных: модуль импорта считывает сжатый бинарный файл выгрузки геоданных, извлекая сырые вершины, ребра и отношения. Данные загружаются во временную реляционную базу данных с пространственным расширением.
    \item Пространственная фильтрация: выполняется удаление тупиковых ребер, пешеходных и специализированных путей, незамкнутых островных подграфов, нарушающих свойство сильной связности графа. Выполняется привязка координат точек пересечений.
    \item Реберная трансформация графа: реляционный граф конвертируется в направленный реберный граф. Дорожные сегменты преобразуются в узлы вычислительного графа, а физически возможные маневры на перекрестках (с учетом знаков приоритета и запретов поворотов) — в направленные ребра.
    \item Расчет весов и оптимизация кэша: производится расчет статического времени преодоления свободных сегментов T\_free. Узлы переупорядочиваются методом обратного алгоритма Катхилла–Макки (RCM) для обеспечения максимальной локальности в памяти.
    \item Выбор ориентиров: запускается алгоритм выборки опорных маяков методом максимального покрытия. Выполняется расчет матрицы расстояний от всех узлов до маяков.
    \item Формирование структуры графа и сериализация: топология графа пакуется в плоские бинарные массивы компактного формата, выравнивается по размеру кэш-линии (64 байта) и сериализуется в единый оптимизированный бинарный файл графа, готовый для загрузки в оперативную память вычислительного ядра через системный вызов отображения файлов в память за константное время.
\end{enumerate}

Схематичная структура хранения компактного представления графа в бинарном дампе памяти приведена выше на рис. \ref{fig:binary_dump}.

\section{Алгоритмическая реализация Time-Dependent A* и SIMD-оптимизации}

Сердцем модуля \texttt{router} является класс \texttt{TdAltRouter}, реализующий алгоритм поиска кратчайшего пути с временной зависимостью (Time-Dependent A*). Для обеспечения заявленной пропускной способности в тысячи запросов в секунду код маршрутизатора глубоко оптимизирован с использованием векторных инструкций процессора.

\subsection{Векторизация 4-арной кучи (AVX2)}

В классической реализации алгоритма Дейкстры или A* основным узким местом (занимающим до 60\% процессорного времени) является извлечение узла с минимальным весом из очереди с приоритетом (Priority Queue). Традиционно для этого используется бинарная куча (\texttt{std::priority\_queue}), имеющая глубину $\log_2(V)$.

В разработанном ядре осуществлен переход на \textbf{4-арную кучу} (4-ary Heap), которая уменьшает глубину дерева в два раза ($\log_4(V)$), что критически важно для снижения количества промахов кэша. Однако в 4-арной куче операция проталкивания узла вниз (Sift-down) требует нахождения минимума среди 4 потомков на каждом уровне, что увеличивает количество инструкций сравнения и ветвлений.

Для нивелирования этих затрат в проекте реализован поиск минимального потомка с использованием 256-битных векторных инструкций AVX2 (Advanced Vector Extensions).
Четыре 64-битные записи потомков (содержащие вес и ID узла) загружаются в один YMM-регистр:
\begin{verbatim}
inline int find_min_child_avx2(const uint32_t* data_ptr) {
    __m256i data = _mm256_loadu_si256(reinterpret_cast<const __m256i*>(data_ptr));
    // Наложение маски для извлечения 32-битных весов из 64-битных записей
    __m256i id_mask_and = _mm256_set1_epi64x(0xFFFFFFFF00000000);
    __m256i weights_only = _mm256_and_si256(data, id_mask_and);
    
    // Поиск минимального значения (через перестановки) и генерация маски
    // ...
    int child_idx = __builtin_ctz(bitmask) / 8;
    return child_idx;
}
\end{verbatim}
Векторизация устраняет все условные переходы (\texttt{if-else}) при поиске дочернего узла. Процессор выбирает минимального потомка за несколько аппаратных тактов без риска сброса конвейера из-за ошибочного предсказания переходов (Branch Misprediction).

\subsection{Механизм Lock-Free резервирования и интерполяции (LERP)}

В процессе раскрытия узлов (Expand) алгоритм TD-A* оценивает вес каждого маневра. Для чтения объемов трафика из \texttt{volume\_buckets} маршрутизатор не использует транзакционные мьютексы. Массивы атомарных переменных считываются со слабым порядком памяти (\texttt{memory\_order\_relaxed}). 

Для устранения ступенчатости штрафов реализован механизм линейной интерполяции (LERP) времени, который считывает объемы из двух соседних 5-минутных слотов ($P_1$ и $P_2$) и интерполирует непосредственно квадраты плотности потока, применяя заранее вычисленный \texttt{K\_magic}:
\begin{verbatim}
uint32_t p1 = K_magic * v1 * v1;
uint32_t p2 = K_magic * v2 * v2;
uint32_t bpr_penalty = (p1 + (p2 - p1) * (eta % 300) / 300) >> 20;
\end{verbatim}
После того как оптимальный путь (цепочка сегментов) найден, роутер осуществляет резервирование емкости. Проходя по восстановленному пути, для каждого сегмента вызывается операция \texttt{fetch\_add}, которая атомарно инкрементирует значение в соответствующей временной корзинке. Это математически аппроксимирует эффект перераспределения потока (Congestion Pricing), уводя следующих агентов на альтернативные маршруты.

\subsection{Пространственные индексы МПР (In-Memory SpatialGrid)}

Для обеспечения мгновенной реакции Модуля принятия решений (МПР) на инциденты реализован сверхбыстрый пространственный индекс. Классические деревья (R-Tree), основанные на указателях в куче, были заменены на плоскую регулярную сетку (SpatialGrid), развернутую в непрерывном массиве. 

Привязка географической геометрии дорог к ячейкам сетки выполняется на этапе офлайн-препроцессинга. В реальном времени МПР формирует ограничивающий прямоугольник (Bounding Box) зоны ДТП или перекрытия и за $O(1)$ получает доступ к списку ID ребер, попавших в ячейки. Затем МПР выставляет директивные штрафы в атомарное поле \texttt{mpr\_penalty} целевых сегментов. Благодаря строгой архитектуре кэширования, поиск затронутых ребер и обновление штрафов происходят за доли микросекунды, не блокируя работу соседних потоков пула маршрутизатора.

\section{Подсистема телеметрии и интеграция с внешними клиентами}

Сбор и трансляция метрик загруженности дорог в реальном времени (Heatmap) является ресурсоемкой задачей, способной существенно замедлить работу основного вычислительного конвейера. Для решения этой проблемы в \texttt{traffic\_core} реализован асинхронный модуль телеметрии, интегрированный с шиной сообщений ZeroMQ.

\subsection{Асинхронная архитектура трансляции состояния (ZeroMQ)}

Вместо прямых TCP-соединений или HTTP-опросов со стороны внешних систем, ядро симуляции выступает в роли однонаправленного источника данных (Publisher). Для организации потоковой передачи применена легковесная библиотека ZeroMQ (паттерн Pub-Sub). 

На этапе «Шаг телеметрии» в главном цикле (Main Loop) симулятор (\texttt{data\_provider}) формирует массив только измененных состояний сегментов (дельт). Если плотность потока на ребре изменилась, формируется компактная структура данных, содержащая \texttt{edge\_id} и текущий расчетный штраф BPR (или цветовую кодировку загруженности). 
Сформированные массивы бинарных данных сериализуются (с применением структур, аналогичных Protobuf/FlatBuffers для обеспечения компактности) и атомарно передаются в фоновый поток Publisher-а ZeroMQ через неблокирующую очередь \texttt{ConcurrentQueue}. 

Такая архитектура гарантирует, что если внешний клиент (например, API Gateway или Web-браузер) отключится, работает медленно или не успевает читать сокет, внутренний цикл సిмулятора C++ не заблокируется (Non-blocking I/O). Сообщения накапливаются во внутренних буферах ZeroMQ или сбрасываются при переполнении, сохраняя строгий бюджет времени на симуляцию.

\subsection{Формирование Heatmap и агрегация метрик}

Для визуализации дорожной обстановки (тепловой карты заторов — Heatmap) на клиентской стороне, \texttt{traffic\_core} не занимается рендерингом тайлов. Вычислительное ядро транслирует только сырые факты: идентификатор направленного ребра и уровень его перенасыщения (отношение $V(t) / C_{max}$). 

МПР (Модуль принятия решений) также использует этот телеметрический контур для сбора глобального показателя эффективности системы — Travel Time Index (TTI). Поскольку ядро (\texttt{router}) фиксирует плановое время прибытия (ETA), а \texttt{data\_provider} отслеживает фактическое завершение поездки агента, система на лету агрегирует показатель $TTI = W_{real} / T_{free}$ для каждого успешно доехавшего агента. 

Асинхронный экспорт этих метрик через ZeroMQ позволяет внешнему API Gateway (написанному на Python или Go) собирать статистику, сохранять исторические срезы в базу данных (ClickHouse/PostgreSQL) и раздавать агрегированные тепловые карты десяткам веб-клиентов (на базе MapLibre GL JS) по WebSocket без какого-либо дополнительного влияния на ядро маршрутизатора.


\section{Выводы по второй главе}
В рамках второй главы была подробно описана программная реализация высокопроизводительного ядра маршрутизации. Реализована строгая архитектурная изоляция вычислительных слоев на основе паттернов инверсии управления и эффективной организации данных в памяти. 

Разработанный конвейер предварительной обработки позволяет эффективно сжимать реляционные картографические данные из базы данных в плоские бинарные массивы компактного формата, применяя алгоритмы топологической пересортировки для минимизации промахов процессора мимо кэша. Использование векторных SIMD-инструкций для работы с многоарными кучами и пространственными сетками обеспечило кратное ускорение критических участков кода, а механизм отображения файлов в виртуальную память полностью устранил задержки при холодном старте системы. Внедренная шина событийного обмена сообщениями позволила отвязать передачу тяжелой телеметрии, включающей карты плотности заторов, от вычислительного потока, гарантируя стабильную работу ядра в режиме жесткого реального времени.
